\begin{prop}
\label{prop:cexEPT}
There exists a K3 surface $S$ over $R:=\bigcup_n\R\lpr t^{1/n}\rpr$ that does not satisfy the EPT property.
\end{prop}

\begin{proof}
Consider $\P^1_R\times \P^1_R$ with homogeneous coordinates $([x:y],[v:w])$.
Let $T$ be the double cover of $\P^1_R\times\P^1_R$, with branch locus of bidegree $(4,4)$, defined by the equation
\[
z^2=x(y-x)(xw^2-yv^2)(xw^2+yv^2)-tw^4y^4.
\]
Since $T$ has only rational double points as singularities, its minimal resolution of singularities $S$ is a K3 surface over~$R$. We denote by $i:T\to T$ the involution associated with the double cover.

The unique connected component $K$ of $S(R)$ contained in $0<x/y<1$ is semi-algebraically isomorphic to a sphere, hence satisfies $H^1(K,\Z/2\Z)=0$. The strict transform $\bar{D}$ of $D:=\{v=0\}$ in $S$ has real points only in $K$, so that its Borel-Haefliger class is trivial. Suppose for contradiction that $S$ satisfies the EPT property. Then $\bar{D}$ is linearly equivalent to a divisor whose support has no $R$-points. Pushing forward to $T$ shows that $D$ is linearly equivalent in $T$ to a Weil divisor whose support contains only finitely many $R$-points.

Write $D=E_1-E_2+\divi(f)$, where $E_1$ and $E_2$ are effective Weil divisors whose supports contain finitely many $R$-points. The divisor $E_2+i(E_2)$ comes from $\P^1\times \P^1$, hence is the zero locus of a section in $H^0(T,\sO(d_1,d_2))$ for some $d_1, d_2\in\Z$. It follows that $E_1+i(E_2)$ is the zero locus of a section in $H^0(T,\sO(d_1,d_2+1))$. These two divisors have finitely many $R$-points in their support. Consequently, after possibly replacing~$d_2$ by $d_2+1$, we have constructed $s\in H^0(T,\sO(d_1,d_2))$ with $d_2$ odd such that $\{s=0\}$ contains finitely many $R$-points.

Suppose there exists $a\in T(R)$ such that $s(a)\neq 0$, $i^*s(a)\neq 0$ but $s(a)+i^*s(a)=0$. Choose a semi-algebraic path $\gamma:[0,1]\to T(R)$ joining $a$ and a point $b\in T(R)$ in the ramification locus. Since $s$ vanishes at only finitely many $R$-points, it is possible to assume (after changing $b$) that neither $s$ nor $i^*s$ vanishes along the path. Since $b$ is a ramification point, $s(b)=i^*s(b)$. By hypothesis, $s(a)=-i^*s(a)$.
Since $[0,1]$ is semi-algebraically connected,
the continuous semi-algebraic map
\[
\lambda\mto i^*s(\gamma(\lambda))/s(\gamma(\lambda))
\]
must reach the value $0$, a contradiction. We have proven that $s(a)+i^*s(a)=0$ implies that $s(a)=0$ or $i^*s(a)=0$, hence that $s+i^*s$ vanishes at only finitely many $R$-points of $T$.

Since $s+i^*s$ is $i^*$-invariant, it comes from a
section $s'\in H^0(\P^1_R\times\P^1_R,\sO(d_1,d_2))$ that vanishes at only finitely many
$R$\nobreakdash-points outside of the closed semi-algebraic subset
$\Theta \subset \P^1(R)\times \P^1(R)$
defined by the inequality
\begin{equation}
\label{eq:theta inequality}
x(y-x)(xw^2-yv^2)(xw^2+yv^2)\leq tw^4y^4\rlap{.}
\end{equation}
Replacing $B:=\{s'=0\}$ by an appropriate reduced irreducible component, we
may assume that $B$ is integral.

Let $A=\R\lcr t^{1/n}\rcr$, with~$n$ large enough that
$B=\sB\otimes_A R$ for some closed integral
subscheme $\sB \subset \P^1_A \times \P^1_A$.
Let~$\sB'$ be the normalisation of~$\sB$ and $B'=\sB'\otimes_AR$.
The map $\sB'\otimes_A \R \to \P^1_\R$ induced by the first projection
is generically finite of degree~$d_2$, which is odd;
therefore $\sB' \otimes_A \R$ possesses an irreducible component~$Y$ of odd multiplicity~$m$,
which dominates~$\P^1_\R$ with odd degree.
Letting $Y_\rs \subset \sB'_\rs \supset B'_\rs$ denote the real spectra of~$Y$, $\sB'$
and~$B'$,
the inequality~\eqref{eq:theta inequality} now defines a closed subset~$\Theta_A$ of~$\sB'_\rs$
which contains $B'_\rs$.
The completion of the local ring of~$\sB'$ at the generic point of~$Y$ is isomorphic,
over~$A$, to $\R(Y)\lcr u\rcr$, where $u=(\alpha t^{1/n})^{1/m}$ for some $\alpha \in \R(Y)^*$.
As~$m$ is odd,
any ordering of~$\R(Y)$ can be extended to $\R(Y)\lpr u\rpr$, compatibly with the given ordering on~$A$,
by declaring that $u$ is infinitely small of the same sign as~$\alpha$.
Thus, any point of~$Y_\rs$ above the generic point of~$Y$ belongs to the closure of~$B'_\rs$ in~$\sB'_\rs$
(see \cite[Prop.\,7.1.21]{bcr}) and, hence, to~$\Theta_A$.
Letting~$Y'$ be the normalisation of~$Y$, it follows that the image of~$Y'_\rs$ in~$Y_\rs$
is contained in~$\Theta_A$.

We have now constructed a smooth, proper, irreducible curve~$Y'$ over~$\R$, endowed with two rational
functions $x,v \in \R(Y')^*$, such that the map $x:Y'\to \P^1_\R$ has odd degree and such that
the function
\[
g=x(1-x)(x-v^2)(x+v^2)
\]
takes negative values on~$Y'(\R)$
wherever it is invertible.
There exists $P \in Y'(\R)$ at which~$x$ vanishes to odd order.
The function~$g$ then has a zero or pole of odd order at~$P$; its
sign therefore changes around~$P$, a contradiction.
\end{proof}
